% TEMPLATE-MILC-v3.tex - plantilla maestra UNICA de las guias MILC v3.
%
% La usan las tres familias de guias, cada una con su driver:
%   · Grados 6-11  -> scripts/build-guias-g11.py
%   · Bebras       -> scripts/build-guias-bebras.py
%   · Semillero    -> scripts/build-guias-semillero.py
%
% Antes habia tres copias casi identicas (~400 lineas iguales, 6 distintas):
% cada mejora habia que aplicarla tres veces, y en la practica no se hacia
% (el motor de recursos vivia solo en dos de ellas). Lo que varia entre
% familias esta parametrizado en 6 placeholders que cada driver llena
% (se nombran aqui SIN su sintaxis real: el builder reemplaza por texto plano,
% tambien dentro de los comentarios, y un valor multilinea rompe el preambulo):
%
%   HEADER_IZQ        encabezado izquierdo de pagina
%   HEADER_DER        encabezado derecho de pagina
%   PORTADA_ETIQUETA  rotulo sobre el numero grande (GRADO/MOMENTO/MODULO)
%   PORTADA_SERIE     linea de serie de la portada
%   PORTADA_ID        identificador de la guia en la portada
%   PORTADA_CIRCULO   el circulito de grado (°), o vacio si no aplica
%
% El resto de claves las documenta content/guias/_SCHEMA.md. El script valida
% que no queden placeholders sin reemplazar antes de escribir el archivo final.

\documentclass[11pt,letterpaper]{article}
\usepackage[margin=1.75cm]{geometry}
\usepackage[table]{xcolor}
\usepackage{fontspec}
\usepackage[spanish]{babel}
\usepackage{tikz}
% Librerías TikZ para los diagramas de las guías (bloque `recursos:`):
%  · babel  → babel en español vuelve activos caracteres como «>», y TikZ los usa
%             en sus opciones (>=stealth, ->). Sin esta librería, cualquier
%             diagrama con flechas revienta con «\language@active@arg> has an
%             extra }». Debe ir ANTES de que se dibuje nada.
%  · positioning        → sintaxis «below left=2pt and 3pt».
%  · decorations.pathreplacing → llaves ({brace}) para señalar rangos.
\usetikzlibrary{babel,positioning,decorations.pathreplacing}
\usepackage{graphicx}
% Circuitos: el trabajo de electrónica se hace hoy con micro:bit y sus sensores
% integrados (MakeCode, sin cableado), así que no se necesitan esquemáticos.
% Si algún día entran montajes con protoboard, basta descomentar esta línea:
% los diagramas .tex/.tikz declarados en `recursos:` ya se incrustan solos.
% \usepackage{circuitikz}
\usepackage[most]{tcolorbox}
\usepackage{tabularx}
\usepackage{array}
\usepackage{fancyhdr}
\usepackage{titlesec}
\usepackage{ragged2e}
\usepackage{enumitem}
\usepackage{microtype}

% Helvetica Neue no trae la flecha U+2192, asi que cada «→» escrito literalmente
% en el YAML se compilaba a NADA, en silencio (462 flechas en 61 guias: rutas del
% tipo «paleta → tipografia → lockup» salian rotas). Mapeamos el caracter a su
% version matematica, que si tiene glifo. Asi el autor puede seguir escribiendo
% «→» comodamente en el YAML.
\usepackage{newunicodechar}
\newunicodechar{→}{\ensuremath{\rightarrow}}

\setmainfont{Helvetica Neue}
\setsansfont{Helvetica Neue}
\setlength{\parindent}{0pt}
\setlength{\parskip}{6pt}
\hyphenpenalty=200
\exhyphenpenalty=200
\pretolerance=400
\tolerance=2000
\setlength{\emergencystretch}{1em}
\renewcommand{\arraystretch}{1.25}
\setlist{leftmargin=5mm,itemsep=1mm,topsep=1mm}
\newcolumntype{Y}{>{\RaggedRight\arraybackslash}X}

\definecolor{milcMagenta}{HTML}{E6007E}
\definecolor{milcVino}{HTML}{5A0038}
\definecolor{milcTurquesa}{HTML}{0093A5}
\definecolor{milcVerde}{HTML}{58A923}
\definecolor{milcMostaza}{HTML}{E5B400}
\definecolor{milcOcre}{HTML}{B86600}
\definecolor{milcNegro}{HTML}{241816}
\definecolor{milcGris}{HTML}{F7F7F4}
\definecolor{milcLinea}{HTML}{E8E2DD}
\definecolor{milcGradePrimary}{HTML}{FF2D87}
\definecolor{milcGradeDark}{HTML}{9F1B56}
\definecolor{milcGradeSoft}{HTML}{FFE0EE}
\definecolor{milcGradeAccent}{HTML}{CFE7FF}
\definecolor{milcGradeText}{HTML}{FFFFFF}
\color{milcNegro}

\pagestyle{fancy}
\fancyhf{}
\lhead{\small Tecnología e Informática · Grado 10}
\rhead{\small Guía 22 · MILC v3}
\cfoot{\small\thepage}
\renewcommand{\headrulewidth}{0pt}

\newtcolorbox{titlebox}[1]{
  enhanced,colback=#1,colframe=#1,arc=7mm,boxrule=0pt,
  left=7mm,right=7mm,top=3mm,bottom=3mm,
  before skip=8pt,after skip=8pt,
  fontupper=\sffamily\bfseries\Large\color{white}
}

\newtcolorbox{softbox}[3]{
  enhanced,colback=#2,colframe=#1,borderline west={4pt}{0pt}{#1},
  arc=2mm,boxrule=.4pt,left=4mm,right=4mm,top=3mm,bottom=3mm,
  before skip=6pt,after skip=8pt,title={#3},coltitle=white,
  colbacktitle=#1,fonttitle=\sffamily\bfseries,fontupper=\RaggedRight,
  attach boxed title to top left={xshift=3mm,yshift=-2mm},
  boxed title style={arc=2mm, boxrule=0pt}
}

\newtcolorbox{drawbox}[2]{
  enhanced,colback=white,colframe=#1,arc=4mm,boxrule=1pt,
  left=5mm,right=5mm,top=4mm,bottom=4mm,before skip=8pt,after skip=8pt,
  title={#2},coltitle=white,colbacktitle=#1,fonttitle=\sffamily\bfseries,
  fontupper=\RaggedRight,
  attach boxed title to top left={xshift=4mm,yshift=-2mm},
  boxed title style={arc=3mm, boxrule=0pt}
}

\newtcolorbox{infoband}[1]{
  enhanced,colback=#1!8,colframe=#1!50,arc=2mm,boxrule=.5pt,
  left=4mm,right=4mm,top=2mm,bottom=2mm,before skip=4pt,after skip=4pt,
  fontupper=\small\RaggedRight
}

% Puente narrativo entre secciones (contrato editorial Conéctate).
% Da continuidad: "Acabas de X. Ahora vas a Y porque Z."
% Visualmente: banda izquierda en color del grado, texto en itálica pequeña.
\newtcolorbox{puentebox}{
  enhanced,colback=milcGradeSoft,colframe=milcGradeDark,
  borderline west={3pt}{0pt}{milcGradeDark},
  arc=0mm,boxrule=0pt,left=5mm,right=4mm,top=2.5mm,bottom=2.5mm,
  before skip=4pt,after skip=8pt,
  fontupper=\itshape\small\color{milcNegro}\RaggedRight
}
% La flecha va en modo matematico a proposito: Helvetica Neue NO tiene el glifo
% U+2192, asi que \textrightarrow se compilaba a NADA («Missing character: There
% is no → in font Helvetica Neue») y el puente salia sin flecha. En modo
% matematico la flecha viene de la fuente matematica y si se dibuja.
\newcommand{\puente}[1]{%
  \begin{puentebox}%
    \textcolor{milcGradeDark}{$\rightarrow$}\hspace{2mm}#1%
  \end{puentebox}%
}

\newcommand{\linead}[1]{\noindent\color{milcLinea}\rule{#1}{0.5pt}\color{milcNegro}}
\newcommand{\respuesta}[1]{\par\vspace{2mm}\foreach \n in {1,...,#1}{\noindent\color{milcLinea}\rule{\linewidth}{0.5pt}\par\vspace{5mm}}\color{milcNegro}}
\newcommand{\checkbox}{\textcolor{milcGradeDark}{\fbox{\phantom{x}}}\hspace{5pt}}

\newcommand{\phasecard}[4]{
\begin{tcolorbox}[enhanced,colback=#3,colframe=#2,arc=3mm,boxrule=.5pt,height=3.05cm,valign=center,halign=center,left=1.5mm,right=1.5mm,top=2mm,bottom=2mm]
{\sffamily\bfseries\fontsize{22}{24}\selectfont\textcolor{#2}{#1}}\par
{\sffamily\bfseries\small #4}
\end{tcolorbox}
}

% ── Recursos visuales de la guía ──────────────────────────────────────
% Declarados en el bloque `recursos:` del YAML y emitidos por el builder.
% \guiaFigura   → imagen rasterizada o PDF (foto, carta celeste, montaje).
% \guiaDiagrama → fuente TikZ/circuitikz (.tex) incrustada como vector:
%                 es la vía para circuitos y esquemáticos editables.
% #1 = ruta relativa al .tex · #2 = pie (puede ir vacío)
\newcommand{\guiaFigura}[2]{%
\begin{center}
\includegraphics[width=0.88\linewidth,height=0.38\textheight,keepaspectratio]{#1}\\[1.5mm]
{\footnotesize\itshape\textcolor{milcNegro!75}{#2}}
\end{center}
}

\newcommand{\guiaDiagrama}[2]{%
\begin{center}
\input{#1}\\[1.5mm]
{\footnotesize\itshape\textcolor{milcNegro!75}{#2}}
\end{center}
}

\begin{document}
\thispagestyle{empty}
\begin{tikzpicture}[remember picture,overlay]
  \fill[white] (current page.south west) rectangle (current page.north east);
  \path[fill=milcGradePrimary,rounded corners=16mm]
    ([xshift=.55cm,yshift=-.55cm]current page.north west) rectangle
    ([xshift=-.55cm,yshift=3.65cm]current page.south east);

  \node[anchor=north west,text=milcGradeText] at ([xshift=2.0cm,yshift=-1.85cm]current page.north west)
    {\sffamily\bfseries\fontsize{14}{16}\selectfont GRADO};
  \node[anchor=north west,text=milcGradeText,opacity=.82] at ([xshift=2.0cm,yshift=-2.75cm]current page.north west)
    {\sffamily\bfseries\fontsize{9}{11}\selectfont SERIE GUÍAS · TECNOLOGÍA E INFORMÁTICA};

  \begin{scope}[shift={([xshift=-3.05cm,yshift=-2.55cm]current page.north east)}]
    \fill[white,opacity=.80] (-.62,-.52) rectangle (.62,.52);
    \draw[milcGradeText,line width=1pt] (-.62,-.52) rectangle (.62,.52);
    \fill[milcGradeDark] (-.42,-.38) rectangle (-.22,.06);
    \fill[milcGradeAccent] (-.10,-.38) rectangle (.10,.24);
    \fill[milcGradePrimary!60!black] (.22,-.38) rectangle (.42,.38);
  \end{scope}

  \node[anchor=west,text=milcGradeText] at ([xshift=1.9cm,yshift=-12.85cm]current page.north west)
    {\sffamily\bfseries\fontsize{124}{124}\selectfont 10};
  % El circulito de grado (°) solo aplica a las guias de grado («11°»); en
  % Bebras y Semillero el numero grande es un momento/modulo, no un grado.
  \draw[milcGradeText,line width=1.7pt]
    ([xshift=4.52cm,yshift=-11.68cm]current page.north west) circle (.13cm);

  \node[anchor=west,text=milcGradeText,text width=8.7cm,align=left] at ([xshift=10.35cm,yshift=-11.6cm]current page.north west)
    {
     {\sffamily\bfseries\fontsize{19}{22}\selectfont Contabilidad básica ---\\ingresos, egresos\\y balance}\\[4mm]
     {\sffamily\bfseries\fontsize{23}{26}\selectfont Guía 22-10-TIC}\\[2mm]
     {\sffamily\fontsize{13}{16}\selectfont Período 3 · Ofimática con IA: contabilidad, datos y emprendimiento}\\[5mm]
     {\sffamily\bfseries\fontsize{11}{13}\selectfont Institución Educativa Sor María Juliana}\\[1mm]
     {\sffamily\fontsize{10}{12}\selectfont Autor: Dr. Álvaro Cárdenas Orozco}
    };

  \path[fill=milcGradeSoft,rounded corners=7mm]
    ([xshift=1.35cm,yshift=.85cm]current page.south west) rectangle
    ([xshift=-1.35cm,yshift=4.25cm]current page.south east);
  \draw[milcGradeDark,line width=.65pt,rounded corners=7mm]
    ([xshift=1.35cm,yshift=.85cm]current page.south west) rectangle
    ([xshift=-1.35cm,yshift=4.25cm]current page.south east);
  \node[anchor=center,text=milcNegro,text width=17.0cm,align=center] at ([yshift=2.55cm]current page.south)
    {\sffamily\fontsize{10}{12}\selectfont Metodología MILC · Apertura ancestral · Pensamiento computacional · Triángulo Dussel-Estoicismo-Floridi\\
    \textcolor{milcGradeDark}{\bfseries Producto:} Cuaderno de cuentas personal o familiar real de 2 semanas con (a) tabla de ingresos por fuente (mesada, propinas, ventas, regalos) registrados día por día, (b) tabla de egresos por categoría (transporte, alimentación, ocio, ahorro, otros) registrados día por día, (c) cálculo del balance neto (ingresos menos egresos) semanal y total, (d) identificación de al menos 1 categoría donde podrías reducir gastos sin sacrificar lo importante + reflexión personal de 5 líneas sobre lo que descubriste de tus hábitos económicos};
\end{tikzpicture}
\mbox{}
\newpage

\section*{Apertura: Contabilidad básica --- ingresos, egresos y balance personal o familiar}

\begin{softbox}{milcGradeDark}{milcGradeSoft}{Saber ancestral}
En las tiendas del centro de Cartago, en las panaderías del barrio Obrero, en las modistas del Quindío, hubo un \textbf{ritual nocturno} que cualquier microempresario practicaba sin falta: \textbf{la cuenta de la caja}. Cuando cerraban al final del día, no se iban a casa directo. Sacaban cuaderno cuadriculado, calculadora, los billetes y monedas, y hacían \textbf{3 operaciones del oficio}: (1) \textbf{Sumar lo que entró}: todas las ventas del día. (2) \textbf{Restar lo que salió}: todos los gastos. Mercancía, servicios, recursos. (3) \textbf{Calcular el saldo}: ingresos menos egresos. Si era positivo, el negocio iba bien. Si era negativo, había problema. Si se repetía el negativo, el negocio estaba en crisis. Esa contabilidad no era trámite: era \textbf{información que dirigía las decisiones del día siguiente}. Si los ingresos bajaban, había que cambiar algo. Si los egresos subían, había que revisar gastos. La sabiduría era precisa: \textbf{ningún negocio sostenible se llevaba solo con la memoria}. La contabilidad personal moderna hereda esa práctica del tendero.
\end{softbox}

\begin{softbox}{milcTurquesa}{milcGris}{Saber contemporáneo}
La \textbf{contabilidad básica} es la práctica de registrar sistemáticamente entradas y salidas de dinero para entender la salud económica de una persona, familia o negocio. Sus \textbf{3 elementos obligatorios} son: (1) \textbf{Ingresos}: todo dinero que entra. Para un joven: mesada, regalos, propinas, pequeños trabajos, ventas casuales. Para una familia: salarios, ingresos de microemprendimiento, rentas. (2) \textbf{Egresos}: todo dinero que sale. Categorías típicas: transporte, alimentación, ocio, ahorro, salud, vivienda, educación, otros. (3) \textbf{Balance neto}: ingresos menos egresos en un periodo (día, semana, mes). Si es positivo: ganancia o excedente. Si es negativo: déficit o pérdida. La regla profesional: \textbf{si no registras tus gastos, no los controlas; lo que no se mide, no se mejora}. El cuaderno de cuentas puede ser físico (cuaderno cuadriculado tradicional) o digital (Google Sheets, Excel, app móvil como Money Manager, Spendee, GnuCash). Lo importante es la \textbf{disciplina del registro diario}, no la herramienta. Empezar con cuaderno físico durante 2 semanas y después migrar a Sheets es estrategia común.
\end{softbox}

\begin{softbox}{milcMostaza}{milcGris}{Pregunta-puente}
\textit{¿Qué sabía el tendero al hacer la cuenta de la caja cada noche, que el novato olvida cuando lleva sus finanzas \emph{``de memoria''}? ¿Y por qué llevar contabilidad personal a los 16 años puede ahorrarte décadas de problemas financieros?}
\end{softbox}

\begin{softbox}{milcVerde}{milcGris}{Saber y saber hacer}
Al terminar podrás: (1) \textbf{identificar} tus fuentes de ingresos y categorías de egresos personales o familiares; (2) \textbf{aplicar} el registro disciplinado de ingresos y egresos durante al menos 2 semanas; (3) \textbf{analizar} tus datos para identificar patrones (días de mayor gasto, categorías que dominan, momentos de déficit); (4) \textbf{evaluar} dónde podrías reducir gastos sin sacrificar lo importante.
\end{softbox}



\small
\begin{tabularx}{\textwidth}{>{\bfseries}p{3.0cm}Y}
\rowcolor{milcGradePrimary!12}Autor & Dr. Álvaro Cárdenas Orozco\\
DBA & Aplicación de ofimática asistida por IA a contabilidad básica y emprendimiento (MEN, T\&I)\\
Referentes & Lean Startup · Phronesis financiera · Ética del dato empresarial\\
Producto final & Cuaderno de cuentas personal o familiar real de 2 semanas con (a) tabla de ingresos por fuente (mesada, propinas, ventas, regalos) registrados día por día, (b) tabla de egresos por categoría (transporte, alimentación, ocio, ahorro, otros) registrados día por día, (c) cálculo del balance neto (ingresos menos egresos) semanal y total, (d) identificación de al menos 1 categoría donde podrías reducir gastos sin sacrificar lo importante + reflexión personal de 5 líneas sobre lo que descubriste de tus hábitos económicos\\
\end{tabularx}
\normalsize

\puente{Antes de empezar, mira el viaje completo. En la sesión 1 mapeaste las herramientas y tareas de ofimática con IA. Hoy aplicas la primera \textbf{contabilidad básica real}. Las próximas sesiones (Sheets + IA, flujo de caja, visualización, mercado, plan de negocio, comunicación empresarial, proyecto integrador, sustentación) usarán estos datos reales como materia prima.}

\begin{titlebox}{milcGradeDark}
Ruta MILC: así vamos a aprender
\end{titlebox}
\begin{tabularx}{\textwidth}{>{\centering\arraybackslash}X>{\centering\arraybackslash}X>{\centering\arraybackslash}X>{\centering\arraybackslash}X}
\phasecard{1}{milcVerde}{milcGris}{Escucha\\[-1mm]\small Observo, escucho y pregunto.} &
\phasecard{2}{milcTurquesa}{milcGris}{Sistematización\\[-1mm]\small Pensamiento computacional.} &
\phasecard{3}{milcMagenta}{milcGris}{Praxis\\[-1mm]\small Construyo y aplico.} &
\phasecard{4}{milcVino}{milcGris}{Evaluación\\[-1mm]\small Reflexión liberadora.}
\end{tabularx}


\puente{Antes de teorizar, vas a hacer una \emph{lista honesta}: piensa en qué gastos tuviste el día de ayer (transporte, comida, ocio, etc.) y suma a mano. Después piensa en cuánto crees que recibes por semana entre mesada, propinas y otros. Esa estimación grosera es la línea base antes del registro disciplinado.}

\begin{titlebox}{milcVerde}
Fase 1 · Escucha
\end{titlebox}

\begin{softbox}{milcVerde}{milcGris}{Escena para escuchar}
\textbf{Actividad 1 · IDENTIFICA --- Estimación inicial de mis cuentas} (15 min · individual). Sin abrir cuaderno ni app aún, estima a memoria: (1) \textbf{Ingresos}: ¿cuánto recibo por semana entre todo? (mesada + propinas + pequeños trabajos + regalos ocasionales). Estima el promedio. (2) \textbf{Egresos}: piensa en los gastos del día de ayer. Anota cada uno y suma. (3) \textbf{Estimación de hábitos}: ¿en qué crees que se va la mayoría del dinero? Tu respuesta intuitiva. Esa estimación es tu línea base: después la compararás con los datos reales del registro.
\end{softbox}

\checkbox Estimé mis ingresos semanales a memoria.\\
\checkbox Sumé los gastos de ayer a memoria.\\
\checkbox Anoté mi hipótesis sobre dónde se va el dinero.

\begin{infoband}{milcVerde}


\textbf{Lo que va al cuaderno (Actividad 1 · IDENTIFICA).} Título: «Actividad 1 --- Estimación inicial». \textbf{Formato:} 3 fichas (ingresos estimados, gastos ayer, hipótesis de hábito). \textbf{Sabes que terminaste cuando} tienes hipótesis cruda para contrastar con datos reales.
\end{infoband}

\puente{Ya tienes estimación. Ahora vas a entender la \textbf{anatomía de la contabilidad básica personal}: las 3 piezas, las categorías típicas de ingresos y egresos, los 5 errores típicos del registro novato.}

\begin{titlebox}{milcTurquesa}
Fase 2 · Sistematización con pensamiento computacional
\end{titlebox}

\begin{softbox}{milcTurquesa}{milcGris}{Cuatro pilares aplicados al tema}
Tu estimación es la hipótesis. Ahora necesitas la \textbf{anatomía profesional de la contabilidad básica personal}.
\end{softbox}

\small
\begin{tabularx}{\textwidth}{>{\bfseries\RaggedRight\arraybackslash}p{3.6cm}Y}
\rowcolor{milcTurquesa!90}\color{white}Pilar & \color{white}\bfseries Aplicación al tema\\
1. Descomposición & Los \textbf{3 elementos obligatorios} de la contabilidad personal se descomponen así: (1) \textbf{Ingresos por fuente}: tabla con columnas Fecha, Fuente, Monto. Fuentes típicas para joven: mesada (semanal o mensual), propinas (ocasionales), pequeños trabajos (ventas, servicios), regalos (cumpleaños, festividades). (2) \textbf{Egresos por categoría}: tabla con columnas Fecha, Categoría, Descripción, Monto. Categorías típicas: \emph{Transporte} (bus, gasolina), \emph{Alimentación} (cafetería, almuerzo, mecato), \emph{Ocio} (películas, salidas, suscripciones), \emph{Ahorro} (lo que apartas), \emph{Estudios} (útiles, materiales), \emph{Otros}. (3) \textbf{Balance neto}: por día, por semana, total. Fórmula: \texttt{balance = suma de ingresos menos suma de egresos}.\\
2. Reconocimiento de patrones & Las \textbf{6 categorías estándar} para egresos personales de joven colombiano: (a) \textbf{Transporte} (bus, gasolina, mototaxi). (b) \textbf{Alimentación} (tienda escolar, almuerzos fuera, mecato). (c) \textbf{Ocio} (cine, salidas, juegos, suscripciones streaming). (d) \textbf{Estudios} (útiles, materiales, fotocopias). (e) \textbf{Ahorro} (lo que apartas conscientemente). (f) \textbf{Otros} (regalos para otros, gastos médicos, etc.). La regla profesional: \textbf{usar las mismas categorías consistentemente}. Cambiar categorías cada semana hace el análisis imposible.\\
3. Abstracción & Lo esencial cabe en una pregunta: \emph{¿lo que crees que pasa con tu dinero coincide con lo que realmente pasa?}. La mayoría de personas, cuando empiezan a registrar, descubren con sorpresa que el ocio o la alimentación informal son mucho más altos de lo que creían. Ese descubrimiento es el valor real de la contabilidad básica: \textbf{convierte la intuición en información}.\\
4. Algoritmo conceptual & Algoritmo: (1) elegir herramienta de registro (cuaderno cuadriculado físico, hoja de Google Sheets, app móvil). (2) crear 2 tablas (ingresos y egresos) con las columnas correctas. (3) registrar cada día durante 2 semanas (todo: hasta los 500 pesos del dulce). (4) cada domingo, sumar ingresos y egresos de la semana y calcular balance. (5) al terminar las 2 semanas, analizar patrones. (6) identificar 1 categoría reducible.\\
\end{tabularx}
\normalsize

\begin{drawbox}{milcTurquesa}{Anatomía de la contabilidad personal}
\textbf{Ingresos:} fecha + fuente + monto. \\[1mm] \textbf{Egresos:} fecha + categoría + descripción + monto. \\[1mm] \textbf{6 categorías:} transporte / alimentación / ocio / estudios / ahorro / otros. \\[1mm] \textbf{Balance:} ingresos menos egresos por periodo. \\[1mm] \textbf{Disciplina:} registro diario sin saltarse días.
\end{drawbox}

\begin{softbox}{milcMostaza}{milcGris}{Errores comunes y su corrección}
(1) \textbf{Saltarse días}: el cuaderno con huecos no permite análisis. (2) \textbf{Olvidar gastos pequeños}: los 500 pesos del dulce suman al final del mes. (3) \textbf{Mezclar contabilidad propia con familiar}: confunde el análisis. (4) \textbf{Cambiar categorías sin avisar}: imposibilita comparación. (5) \textbf{Hacer todo en la cabeza}: si no escribes, no controlas.
\end{softbox}

\begin{infoband}{milcTurquesa}


\textbf{Lo que va al cuaderno (Actividad 2 · APLICA).} Título: «Actividad 2 --- Anatomía contabilidad». \textbf{Formato:} (a) plantilla con 6 categorías; (b) regla de disciplina diaria; (c) los 5 errores con marca. \textbf{Sabes que terminaste cuando} sabes registrar sin pensar mucho.
\end{infoband}

\puente{Tienes el método. Toca aplicarlo: armas tu cuaderno (físico o digital), registras 2 semanas de ingresos y egresos día a día, calculas balance, identificas categoría a reducir.}

\begin{titlebox}{milcMagenta}
Fase 3 · Praxis con pensamiento computacional
\end{titlebox}

\begin{softbox}{milcMagenta}{milcGris}{Construyendo el producto con los cuatro pilares}
\textbf{Actividad 3 · ANALIZA + EVALÚA --- Cuaderno de 2 semanas + análisis + reflexión} (15 min en sesión + 2 semanas de registro · individual). Hora de aplicar. Armas tu cuaderno, registras 2 semanas, analizas, identificas categoría a reducir, reflexionas.
\end{softbox}

\small
\begin{tabularx}{\textwidth}{>{\bfseries\RaggedRight\arraybackslash}p{3.6cm}Y}
\rowcolor{milcMagenta!90}\color{white}Pilar & \color{white}\bfseries Aplicación al producto\\
Descomposición & Tu cuaderno se construye y se mantiene en 4 pasos: (1) elegir herramienta: cuaderno cuadriculado físico es buen primer paso; al periodo 3-4 puedes migrar a Sheets. (2) crear 2 tablas con encabezados. (3) registrar cada día durante 14 días seguidos: ingresos y egresos sin excepción. (4) cada domingo: sumar semana, calcular balance, anotar observaciones.\\
Patrones & Sigue el patrón \textbf{``registrar antes de juzgar''}: las primeras 2 semanas son solo observación. No te juzgues por gastar en lo que gastas. Solo registra con honestidad. Cuando tengas 14 días completos, recién entonces analiza para encontrar patrones y decidir cambios. La phronesis del oficio consiste en \textbf{convertir hábitos invisibles en información visible antes de intentar cambiarlos}.\\
Abstracción & El \textbf{análisis al final de las 2 semanas} responde 5 preguntas: (a) ¿Cuál fue el ingreso total y promedio semanal? (b) ¿Cuál fue el egreso total y promedio semanal? (c) ¿Cuál fue el balance neto (positivo o negativo)? (d) ¿Qué categoría representó el mayor porcentaje de egresos? (e) ¿Dónde podría reducir sin sacrificar lo importante?\\
Algoritmo de armado & Algoritmo: (1) crear cuaderno. (2) registrar 14 días. (3) calcular totales y balance. (4) analizar categorías. (5) identificar 1 categoría reducible. (6) reflexión de 5 líneas sobre lo descubierto.\\
\end{tabularx}
\normalsize

\begin{drawbox}{milcMagenta}{Checklist al cerrar las 2 semanas}
\checkbox 14 días completos de registro (sin huecos). \\[1mm] \checkbox Ingresos por fuente registrados. \\[1mm] \checkbox Egresos por categoría registrados. \\[1mm] \checkbox Balance neto por semana calculado. \\[1mm] \checkbox Categoría dominante identificada. \\[1mm] \checkbox 1 categoría reducible declarada. \\[1mm] \checkbox Reflexión de 5 líneas sobre hábitos descubiertos.
\end{drawbox}

\begin{softbox}{milcMagenta}{milcGris}{Plantilla de trabajo}
\textbf{Cuaderno.} \textit{Físico o digital, 2 tablas.} \\[1mm] \textbf{Registro.} \textit{14 días sin huecos.} \\[1mm] \textbf{Análisis.} \textit{Ingresos, egresos, balance, categoría dominante.} \\[1mm] \textbf{Decisión.} \textit{Categoría a reducir.} \\[1mm] \textbf{Reflexión.} \textit{5 líneas sobre hábitos descubiertos.}
\end{softbox}

\begin{infoband}{milcMagenta}


\textbf{Lo que va al cuaderno (Actividad 3 · ANALIZA + EVALÚA).} Título: «Actividad 3 --- Mi contabilidad de 2 semanas». \textbf{Formato:} (a) tablas completas; (b) cálculos; (c) reflexión final. \textbf{Sabes que terminaste cuando} descubriste algo sobre tus hábitos que no sabías.
\end{infoband}

\puente{Lo que sigue resume \emph{qué} entregas al final de la semana: cuaderno de 2 semanas + balance + análisis + reflexión. El criterio principal: que descubras algo concreto sobre tus hábitos económicos que no sabías.}

\begin{titlebox}{milcMostaza}
Producto de la sesión
\end{titlebox}
\textbf{Cuaderno 2 semanas + balance + análisis + reflexión}.
\begin{softbox}{milcMostaza}{milcGris}{Criterios de calidad}
(1) 14 días de registro sin huecos. (2) Ingresos y egresos categorizados. (3) Balance neto calculado. (4) Categoría dominante identificada. (5) 1 categoría reducible declarada. (6) Reflexión sobre hábitos descubiertos.
\end{softbox}

\puente{Antes de cerrar, mira la contabilidad desde las cinco dimensiones humanas. Llevar cuenta es respeto por tu propio futuro económico; ignorar las cuentas es entregarse a la deriva financiera.}

\begin{titlebox}{milcVino}
Fase 4 · Evaluación liberadora — cinco dimensiones
\end{titlebox}

\begin{softbox}{milcVino}{milcGris}{Sentido de la evaluación}
Evaluar no es solo revisar una nota. Es reconocer qué aprendí, qué cambió en mí, qué emociones manejé, cómo me relaciono con la ciudad y la comunidad, y qué le dejaré a quien viene detrás.
\end{softbox}

\textbf{Mi reflexión liberadora en cinco dimensiones.} Completa cada frase iniciada:

\small
\begin{tabularx}{\textwidth}{>{\bfseries\RaggedRight\arraybackslash}p{3.4cm}Y>{\RaggedRight\arraybackslash}p{4.6cm}}
\rowcolor{milcVino}\color{white}Dimensión & \color{white}\bfseries Mi respuesta (frame starter) & \color{white}\bfseries Algo que descubrí\\
1. Personal & Lo que más cambió en mí fue \linead{6.5cm} & \linead{4.2cm}\\
2. Emocional & La emoción más fuerte fue \linead{2.5cm} y la regulé \linead{3cm} & \linead{4.2cm}\\
3. Ciudadana & En democracia esto importa porque \linead{6cm} & \linead{4.2cm}\\
4. Local (Cartago) & En mi barrio sería útil para \linead{6.5cm} & \linead{4.2cm}\\
5. Intergeneracional & A mi abuela/abuelo le diría \linead{6.5cm} & \linead{4.2cm}\\
\end{tabularx}
\normalsize

\begin{infoband}{milcVino}
\textbf{Para la próxima sustentación.} Vuelve a esta tabla en seis meses. Verás que las respuestas cambian — no porque lo aprendido cambie, sino porque tú cambias. La evaluación liberadora es una conversación contigo a través del tiempo.
\end{infoband}


\puente{Y para terminar, tres voces sobre el cálculo cotidiano. Dussel sobre la contabilidad personal que libera al microempresario, Séneca sobre la disciplina del registro, Floridi sobre el dato económico personal como información cuidada.}

\begin{titlebox}{milcGradeDark}
Triángulo de pensamiento · cierre filosófico
\end{titlebox}

\begin{softbox}{milcVino}{milcGris}{Dussel · lente del nosotros}
\textit{``La contabilidad personal libera al joven de la deriva financiera; ignorar las cuentas es entregarse al sistema que extrae sin información.''}

Cuando llevas cuenta de tus ingresos y egresos, asumes control de tu vida económica. La filosofía de la liberación pide ese acto: \textbf{conocer tu economía es condición para decidir sobre ella}. Quien no lleva cuenta queda a merced de hábitos automáticos y publicidad. La phronesis financiera del joven consiste en \textbf{empezar el registro temprano, cuando los hábitos aún se pueden moldear}.

\textbf{¿Estoy asumiendo control de mi economía con el registro, o entrego mis decisiones a hábitos automáticos?}
\end{softbox}
\linead{\linewidth}\\[1mm]
\linead{\linewidth}

\begin{softbox}{milcMostaza}{milcGris}{Estoicismo · Séneca · cuidado interior}
\textit{``La disciplina del registro diario es virtud; la negligencia económica es vicio que se descubre tarde.''}

Es tentador saltarse el registro algunos días porque \emph{``no pasó nada importante''}. La sabiduría estoica recomienda lo opuesto: \emph{registra todo, incluso los días tranquilos}. La disciplina del registro continuo es lo único que produce información útil. Quien registra solo cuando se acuerda no controla nada.

\textbf{¿Estoy registrando con disciplina diaria, o salto días esperando que la memoria llene los huecos?}
\end{softbox}
\linead{\linewidth}\\[1mm]
\linead{\linewidth}

\begin{softbox}{milcTurquesa}{milcGris}{Floridi · lente de la infoesfera}
\textit{``Los datos económicos personales bien recolectados son información cuidada que fundamenta decisiones de vida.''}

En la era de la información, llevar cuenta personal es ejercicio de soberanía sobre los propios datos. La ética informacional pide que cuides tus propios números antes de delegarlos a apps o entidades. Tu cuaderno de 2 semanas te entrena en una práctica que escalará a toda tu vida económica: \textbf{controlar tu información para controlar tus decisiones}.

\textbf{¿Estoy cuidando mis datos económicos como información profesional, o solo registro porque me lo pidieron?}
\end{softbox}
\linead{\linewidth}\\[1mm]
\linead{\linewidth}

\begin{titlebox}{milcVino}
Compromiso final
\end{titlebox}

\small
\begin{tabularx}{\textwidth}{>{\RaggedRight\arraybackslash}p{3.0cm}YYY}
\rowcolor{milcVino}\color{white}\bfseries Criterio & \color{white}\bfseries Lo logré & \color{white}\bfseries En proceso & \color{white}\bfseries Necesito apoyo\\
Comprensión del tema & Explico ideas centrales con mis palabras. & Reconozco ideas pero falta precisión. & Necesito repasar conceptos base.\\
Pensamiento computacional & Apliqué los 4 pilares al diseñar mi producto. & Apliqué algunos pilares. & Necesito releer la fase 2.\\
Aplicación y producto & Mi producto cumple los criterios y se puede revisar. & Producto funcional con ajustes pendientes. & Aún en construcción.\\
Reflexión 5 dimensiones & Respondí cada dimensión con honestidad. & Respondí algunas con detalle. & Mis respuestas son muy generales.\\
Triángulo de pensamiento & Conecté las 3 voces con mi trabajo real. & Conecté al menos una. & No relaciono las voces con mi proceso.\\
\end{tabularx}
\normalsize

\textbf{Hoy entendí que:}\\[1mm]
\linead{\linewidth}\\[1mm]
\linead{\linewidth}

\textbf{Mi compromiso para la próxima sesión es:}\\[1mm]
\linead{\linewidth}\\[1mm]
\linead{\linewidth}

\begin{softbox}{milcMostaza}{milcGris}{Cita de despedida · Marco Aurelio}
\textit{``El obstáculo en el camino se vuelve el camino. Lo que estorba al camino se vuelve camino.''}\\[1mm]
Cada error de hoy, bien observado, se convierte en pieza del camino que pisarás mañana.
\end{softbox}

\small
\begin{tabularx}{\textwidth}{>{\bfseries}p{3.6cm}Y}
\rowcolor{milcGradeDark!12}Nombre completo: & \linead{10cm}\\
Fecha: & \linead{4cm} \quad Curso/grupo: \linead{4cm}\\
Mi firma: & \linead{9cm}\\
\end{tabularx}
\normalsize

\begin{infoband}{milcGradeDark}
\textbf{Para la próxima sesión.} Llega con tu cuaderno de 2 semanas completo, balance calculado y análisis. En la sesión 3 vas a migrar tu cuaderno a \textbf{Google Sheets con asistencia de IA}: fórmulas automáticas, gráficos, análisis con \texttt{=GEMINI()}. La contabilidad a mano de hoy es la base; Sheets de mañana es la herramienta de escala.
\end{infoband}

\end{document}
